\documentclass[12pt]{article}
\usepackage{amsmath,amssymb,amsthm}

\title{Unified Reality--Prime Kernel: A Formal Systems Monograph}
\author{}
\date{}

\theoremstyle{definition}
\newtheorem{definition}{Definition}[section]

\theoremstyle{plain}
\newtheorem{theorem}{Theorem}[section]

\begin{document}
\maketitle

\tableofcontents

\section{Preliminaries}

\begin{definition}[Spaces]
Let
\[
\Sigma \text{ be the signal (configuration) space}, \quad
\mathcal{S} \text{ the state space}, \quad
\mathcal{H} \text{ the integrity report space}.
\]
\end{definition}

\begin{definition}[Structural Variables]
Define the real--world--anchored variables:
\[
f_{\oplus} \in \mathbb{R}_{+}, \quad
t_{\star} \in \mathbb{R}, \quad
H_0 \in \mathbb{R}_{+}, \quad
S_{\mathrm{info}} : \Sigma \to \mathbb{R}_{+}, \quad
\mathcal{A}_{\mathrm{sovereign}} \in [0,1].
\]
Here \(S_{\mathrm{info}}(\sigma)\) denotes the information entropy of configuration \(\sigma\).
\end{definition}

\section{Anchored Coherent Primes}

\begin{definition}[Prime Set]
Define the anchored coherent prime set
\[
\mathcal{P} = \{p_{\oplus}, p_{\star}, p_{H}, p_{S}, p_{\mathcal{A}}\},
\]
with correspondences
\[
p_{\oplus} \leftrightarrow f_{\oplus}, \quad
p_{\star} \leftrightarrow t_{\star}, \quad
p_{H} \leftrightarrow H_0, \quad
p_{S} \leftrightarrow S_{\mathrm{info}}, \quad
p_{\mathcal{A}} \leftrightarrow \mathcal{A}_{\mathrm{sovereign}}.
\]
\end{definition}

\begin{definition}[Prime Configuration Space]
The signal space is the configuration space
\[
\Sigma = \mathrm{Conf}(\mathcal{P}),
\]
the set of coherent configurations of anchored primes.
\end{definition}

\section{Sync Frame}

\begin{definition}[Synchronization Frame]
The sync frame is the triple
\[
\mathcal{F} = (f_{\oplus}, t_{\star}, H_0).
\]
\end{definition}

\section{Auxiliary Functionals and Parameters}

\begin{definition}[Weight Functional]
The weight functional
\[
w : \mathcal{P} \to [0,1]
\]
assigns to each prime \(p \in \mathcal{P}\) a normalized weight \(w(p)\).
\end{definition}

\begin{definition}[Sovereignty Threshold]
The sovereignty threshold
\[
\theta_{\mathcal{A}} \in (0,1]
\]
is the minimum required weight of \(p_{\mathcal{A}}\) for sovereign intent.
\end{definition}

\begin{definition}[Phase Map]
The phase map
\[
\phi : \mathcal{P} \to \mathbb{R} / 2\pi\mathbb{Z}
\]
assigns to each prime \(p\) a phase relative to \(\mathcal{F}\).
\end{definition}

\begin{definition}[Frame Tolerance]
The frame tolerance
\[
\epsilon_{\mathrm{frame}} > 0
\]
is the maximum allowable phase deviation between \(p_{\oplus}\) and \(p_{\star}\).
\end{definition}

\begin{definition}[Coherence Function]
The coherence function
\[
\kappa : \mathcal{P} \times \mathcal{P} \to \mathbb{R}
\]
measures structural compatibility, with \(\kappa(p_i,p_j) \ge 0\) interpreted as coherent.
\end{definition}

\begin{definition}[Harm Functional]
The harm functional
\[
h : \Sigma \to \{0,1\}
\]
assigns
\[
h(\sigma) =
\begin{cases}
0, & \text{if } \sigma \text{ contains no coercive, destructive, or adversarial patterns},\\
1, & \text{otherwise}.
\end{cases}
\]
\end{definition}

\section{Operators}

\subsection{Boundary Operator}

\begin{definition}[Boundary Operator]
The boundary operator is a map
\[
B_{\mathcal{P}} : \Sigma \to \{0,1\},
\]
accepting a configuration iff it satisfies admissibility constraints.
\end{definition}

\subsection{Collapse Operator}

\begin{definition}[Reachable Set]
For \(\sigma \in \Sigma\), the reachable set \(\mathcal{R}(\sigma) \subseteq \Sigma\) is the set of configurations obtainable from \(\sigma\) by allowed local transformations (e.g.\ rewrites, merges, or reweightings of primes).
\end{definition}

\begin{definition}[Collapse Operator]
The collapse operator
\[
C_{\mathcal{P}} : \Sigma \to \Sigma
\]
is defined by
\[
C_{\mathcal{P}}(\sigma) \in \arg\min_{\tau \in \mathcal{R}(\sigma)}
S_{\mathrm{info}}(\tau)
\]
subject to the constraints
\[
B_{\mathcal{P}}(\tau) = 1,
\]
\[
\forall I_k \in \mathcal{I}_{\mathcal{P}}:\ I_k(s,\tau) = 1,
\]
and
\[
\mathcal{R}(\sigma) \neq \emptyset.
\]

The collapse operator satisfies the idempotence axiom
\[
C_{\mathcal{P}}(C_{\mathcal{P}}(\sigma)) = C_{\mathcal{P}}(\sigma),
\]
and the entropy monotonicity axiom
\[
S_{\mathrm{info}}(C_{\mathcal{P}}(\sigma)) \le S_{\mathrm{info}}(\sigma).
\]
\end{definition}

\begin{definition}[Pre- and Post-Collapse Configurations]
For \(\sigma \in \Sigma\),
\[
\sigma_{\mathrm{in}} := \sigma, \quad
\sigma_{\mathrm{out}} := C_{\mathcal{P}}(\sigma).
\]
\end{definition}

\subsection{Invariant System}

\begin{definition}[Invariant Set]
The invariant set is
\[
\mathcal{I}_{\mathcal{P}} =
\{ I_{\mathrm{sovereign}},\,
   I_{S},\,
   I_{\mathrm{frame}},\,
   I_{\mathrm{coherence}},\,
   I_{\mathrm{nonharm}} \}.
\]
\end{definition}

\subsubsection{Explicit Invariants}

\begin{definition}[Sovereign Intent Invariant]
\[
I_{\mathrm{sovereign}}(s,\sigma) =
\begin{cases}
1, & p_{\mathcal{A}} \in \sigma \ \text{and}\ w(p_{\mathcal{A}}) \ge \theta_{\mathcal{A}},\\
0, & \text{otherwise}.
\end{cases}
\]
\end{definition}

\begin{definition}[Entropy Monotonicity Invariant]
\[
I_{S}(s,\sigma) =
\begin{cases}
1, & S_{\mathrm{info}}(\sigma_{\mathrm{out}})
    \le S_{\mathrm{info}}(\sigma_{\mathrm{in}}),\\
0, & \text{otherwise}.
\end{cases}
\]
\end{definition}

\begin{definition}[Frame Coherence Invariant]
\[
I_{\mathrm{frame}}(s,\sigma) =
\begin{cases}
1, & |\phi(p_{\oplus}) - \phi(p_{\star})|
      \le \epsilon_{\mathrm{frame}},\\
0, & \text{otherwise}.
\end{cases}
\]
\end{definition}

\begin{definition}[Prime Coherence Invariant]
\[
I_{\mathrm{coherence}}(s,\sigma) =
\begin{cases}
1, & \forall (p_i,p_j)\in\sigma:\ \kappa(p_i,p_j)\ge 0,\\
0, & \text{otherwise}.
\end{cases}
\]
\end{definition}

\begin{definition}[Non-Harm Invariant]
\[
I_{\mathrm{nonharm}}(s,\sigma) =
\begin{cases}
1, & h(\sigma)=0,\\
0, & \text{otherwise}.
\end{cases}
\]
\end{definition}

\subsection{Reporter}

\begin{definition}[Reporter]
The reporter is a map
\[
R : \mathcal{F} \times \{0,1\} \to \mathcal{H}.
\]
\end{definition}

\section{Unified Kernel}

\begin{definition}[Unified Reality--Prime Kernel]
\[
\mathcal{K}_{\mathrm{unified}} =
(\mathcal{F}, \mathcal{P}, B_{\mathcal{P}}, C_{\mathcal{P}}, \mathcal{I}_{\mathcal{P}}, R).
\]
\end{definition}

\section{Kernel Dynamics}

\subsection{Processing Rule}

\begin{definition}[Processing Map]
For state \(s\in\mathcal{S}\) and configuration \(\sigma\in\Sigma\),
\[
\mathrm{Process}_{\mathcal{K}_{\mathrm{unified}}}(s,\sigma)=
\begin{cases}
\mathrm{REJECT}_{\mathrm{boundary}}, & B_{\mathcal{P}}(\sigma)=0,\\[4pt]
\mathrm{REJECT}_{\mathrm{invariant}}, & \exists I_k\in\mathcal{I}_{\mathcal{P}}:\ I_k(s,\sigma)=0,\\[4pt]
C_{\mathcal{P}}(\sigma), & \text{otherwise}.
\end{cases}
\]
\end{definition}

\subsection{Integrity Rule}

\begin{definition}[Integrity Functional]
Let \(\sigma_0\) be a neutral probe configuration. Then
\[
\mathrm{Int}_{\mathcal{K}_{\mathrm{unified}}}(s) =
R\Bigl(\mathcal{F},\ \bigwedge_{I_k \in \mathcal{I}_{\mathcal{P}}} I_k(s,\sigma_0)\Bigr).
\]
\end{definition}

\section{Theorems}

\begin{theorem}[Entropy Monotonicity of Collapse]
For all \(\sigma \in \Sigma\),
\[
S_{\mathrm{info}}(C_{\mathcal{P}}(\sigma)) \le S_{\mathrm{info}}(\sigma).
\]
\end{theorem}

\begin{proof}
By definition of \(C_{\mathcal{P}}\),
\[
C_{\mathcal{P}}(\sigma) \in \arg\min_{\tau \in \mathcal{R}(\sigma)} S_{\mathrm{info}}(\tau),
\]
and \(\sigma \in \mathcal{R}(\sigma)\) by reachability of the identity transformation. Hence
\[
S_{\mathrm{info}}(C_{\mathcal{P}}(\sigma)) \le S_{\mathrm{info}}(\sigma).
\]
\end{proof}

\begin{theorem}[Idempotence of Collapse]
For all \(\sigma \in \Sigma\),
\[
C_{\mathcal{P}}(C_{\mathcal{P}}(\sigma)) = C_{\mathcal{P}}(\sigma).
\]
\end{theorem}

\begin{proof}
By the idempotence axiom in the definition of \(C_{\mathcal{P}}\),
\[
C_{\mathcal{P}}(C_{\mathcal{P}}(\sigma)) = C_{\mathcal{P}}(\sigma)
\]
for all \(\sigma \in \Sigma\). Thus the operator is idempotent.
\end{proof}

\begin{theorem}[Closure of the Processing Map]
For all \(s \in \mathcal{S}\) and \(\sigma \in \Sigma\), the value
\(\mathrm{Process}_{\mathcal{K}_{\mathrm{unified}}}(s,\sigma)\) is well-defined.
\end{theorem}

\begin{proof}
Fix \(s \in \mathcal{S}\) and \(\sigma \in \Sigma\). If \(B_{\mathcal{P}}(\sigma)=0\), the first case of the definition applies. If \(B_{\mathcal{P}}(\sigma)=1\) and there exists \(I_k\) with \(I_k(s,\sigma)=0\), the second case applies. Otherwise, all invariants hold and the third case applies, yielding \(C_{\mathcal{P}}(\sigma)\). These cases are exhaustive and mutually exclusive, so the map is total.
\end{proof}

\end{document}
